% GENERATED FILE. Edit canonical lessons, then run npm run generate:books.
% canonical-source-hash: fnv1a64:4645f0df312ea78d
% canonical-lessons: ZH-C297-baoxian, ZH-C297-shigu, ZH-W297-gu-old, ZH-C297-zeren, ZH-C297-houguo, ZH-C297-chuli

\chapter{Whose Fault}
\label{ch:zh-whose-fault}
\hlchaptermodality{297}

\begin{chapteropening}
\textbf{By the end of this chapter:} \emph{I can say who was to blame for an accident, what came of it and who dealt with it.}

Complete the last lesson of chapter 297: I can say who was to blame for an accident, what came of it and who dealt with it.
\end{chapteropening}

\section[bǎoxiǎn]{bǎoxiǎn — insurance}

\label{lesson:ZH-C297-baoxian}



\begin{quote}
Before the new one: say the Mandarin for a beginner at the wheel, then the Mandarin for to pay for damage.
\end{quote}

\subsection*{You'll want to know: bǎoxiǎn}
\emph{bǎoxiǎn} — \textquotedblleft{}insurance\textquotedblright{}.

\begin{quote}
\emph{wǒ de chē yǒu bǎoxiǎn} — My car is insured.
\end{quote}

Its characters come later; for now it is a word for the ear.

Two third tones meet in this word: bǎo is said with a rising tone before xiǎn, and the spelling keeps both marks.

\begin{cousinweb}
\emph{bǎoxiǎn} is built from \emph{bǎo} 'to guard' and \emph{xiǎn} 'danger', both inherited from Middle Chinese, \emph{pawX} and \emph{xjemX} (Guangyun, in Baxter's transcription), which glosses the first 'to guard' and 'safe' and the second 'danger': a guard against danger.
\end{cousinweb}

\begin{grammarlens}[title={What you've built}]
One more word for telling what happened.
\end{grammarlens}

\subsection*{Your turn}
\begin{itemize}
\raggedright
  \item \emph{Say it:} \emph{bǎoxiǎn}
  \item \emph{Say it:} \emph{bǎoxiǎn}, once more
  \item \emph{Say it:} \emph{wǒ de chē yǒu bǎoxiǎn} — My car is insured.
  \item \emph{Recall:} say the Mandarin for a beginner at the wheel, then the Mandarin for to pay for damage, then say \emph{bǎoxiǎn} again
\end{itemize}

\begin{tcolorbox}[breakable,colback=teal!4,colframe=teal!35!black,title={Before you move on}]
What does bǎoxiǎn mean? (\textquotedblleft{}Insurance\textquotedblright{}.) Say it once more.
\end{tcolorbox}
\section[\texorpdfstring{\zh{事故} (shìgù)}{事故 (shìgù)}]{\zh{事故} (shìgù) — accident}

\label{lesson:ZH-C297-shigu}



\begin{quote}
Before the new one: say the Mandarin for to pay for damage, then the Mandarin for insurance for the car.
\end{quote}

\subsection*{You'll want to know: \zh{事故}}
\textbf{\zh{事故}} — \emph{shìgù} — \textquotedblleft{}accident\textquotedblright{}.

\begin{quote}
\textbf{\zh{这个路口常常发生事故}} \emph{zhège lùkǒu chángcháng fāshēng shìgù} — Accidents often happen at this crossing.
\end{quote}

The new character, in pieces: \zh{故} — nine strokes in two pieces, the right two pieces of \zh{做} without its person piece: on the left \zh{十} over \zh{口}, five strokes; on the right the striking piece, four strokes. The next lesson puts it in your hand.

\begin{cousinweb}
\textbf{\zh{事故}} (\emph{shìgù}) is built from \zh{事} 'a matter' and \zh{故}, both inherited from Middle Chinese, \emph{dzriH} and \emph{kuH} (Guangyun, in Baxter's transcription), which glosses the second 'old' and 'a matter'; how the pair came to mean an accident is not settled.
\end{cousinweb}

\begin{grammarlens}[title={What you've built}]
One more word for telling what happened.
\end{grammarlens}

\subsection*{Your turn}
\begin{itemize}
\raggedright
  \item \emph{Say it:} \emph{shìgù}
  \item \emph{Say it:} \emph{shìgù}, once more
  \item \emph{Say it:} \emph{zhège lùkǒu chángcháng fāshēng shìgù} — Accidents often happen at this crossing.
  \item \emph{Recall:} say the Mandarin for to pay for damage, then the Mandarin for insurance for the car, then say \emph{shìgù} again
\end{itemize}

\begin{tcolorbox}[breakable,colback=teal!4,colframe=teal!35!black,title={Before you move on}]
What does \zh{事故} mean? (\textquotedblleft{}Accident\textquotedblright{}.) Say it once more.
\end{tcolorbox}
\section[\texorpdfstring{\zh{故} (gù)}{故 (gù)}]{\zh{故} — \zh{十} over \zh{口}, then the striking piece}

\label{lesson:ZH-W297-gu-old}



\begin{quote}
The word first, then the character inside it.

\begin{itemize}
\raggedright
  \item \emph{Recall:} say \emph{shìgù} — \textbf{R1}, one lesson back, before it decays
\end{itemize}

\begin{quote}
\zh{做} \zh{十}
\end{quote}

\emph{Read it:} each one aloud, with nothing said first — \emph{zuò}, \emph{shí}
\end{quote}

\subsection*{Script — \zh{十} over \zh{口}, then the striking piece}
\hlblockfigure{\detokenize{figures/ZH-W297-gu-old-filmstrip.pdf}}{How \zh{故} is written, stroke by stroke}

\begin{quote}
\zh{故}
\end{quote}

\textbf{Two pieces, left before right.} Count the pieces, not the strokes.

\textbf{The left piece, five strokes.} \zh{十}: a bar and a vertical through it. Under it \zh{口}: a short vertical, a stroke that runs right and turns down, and the bar that closes the foot.

\textbf{The right piece, four strokes.} The striking piece of \zh{做}: a sweep down to the left, a bar across to the right, a long sweep down to the lower left, and a long stroke falling to the right across it.

5 + 4 = nine.

\subsection*{Your turn}
\begin{itemize}
\raggedright
  \item \emph{Write it:} \zh{故}, all nine, with the model in front of you
  \item \emph{Cover:} write it again from memory, and count the pieces, not the strokes
\end{itemize}

\begin{tcolorbox}[breakable,colback=teal!4,colframe=teal!35!black,title={Before you move on}]
How many pieces in \zh{故}? (\textbf{Two}.) Which character has the same two pieces on its right? (\textbf{\zh{做}}.) What is on the left? (\textbf{\zh{十} over \zh{口}}.)

Source: \href{https://raw.githubusercontent.com/chanind/hanzi-writer-data/68d10a4b21150cae5e1ebbd223eed289cf32d90c/data/\%E6\%95\%85.json}{Hanzi Writer Data \zh{故}, pinned stroke paths and medians}.
\end{tcolorbox}
\section[zérèn]{zérèn — responsibility; blame}

\label{lesson:ZH-C297-zeren}



\begin{quote}
Before the new one: say the Mandarin for insurance for the car, then the Mandarin for a serious accident.
\end{quote}

\subsection*{You'll want to know: zérèn}
\emph{zérèn} — \textquotedblleft{}responsibility; blame\textquotedblright{}.

\begin{quote}
\emph{shìgù fāshēng yǐhòu, sījī shuō shì tā de zérèn} — After the accident, the driver said it was his fault.
\end{quote}

Its characters come later; for now it is a word for the ear.

\begin{cousinweb}
\emph{zérèn} is built from \emph{zé} 'to demand' and \emph{rèn} 'to bear', both inherited from Middle Chinese, \emph{tsreak} and \emph{nyim} (Guangyun, in Baxter's transcription), which glosses the first 'to seek' and the second 'able to bear', and lists the second also in the departing tone, the source of rèn.
\end{cousinweb}

\begin{grammarlens}[title={What you've built}]
One more word for telling what happened.
\end{grammarlens}

\subsection*{Your turn}
\begin{itemize}
\raggedright
  \item \emph{Say it:} \emph{zérèn}
  \item \emph{Say it:} \emph{zérèn}, once more
  \item \emph{Say it:} \emph{shìgù fāshēng yǐhòu, sījī shuō shì tā de zérèn} — After the accident, the driver said it was his fault.
  \item \emph{Recall:} say the Mandarin for insurance for the car, then the Mandarin for a serious accident, then say \emph{zérèn} again
\end{itemize}

\begin{tcolorbox}[breakable,colback=teal!4,colframe=teal!35!black,title={Before you move on}]
What does zérèn mean? (\textquotedblleft{}Responsibility; blame\textquotedblright{}.) Say it once more.
\end{tcolorbox}
\section[\texorpdfstring{\zh{后果} (hòuguǒ)}{后果 (hòuguǒ)}]{\zh{后果} (hòuguǒ) — consequences; what comes of it}

\label{lesson:ZH-C297-houguo}



\begin{quote}
Before the new one: say the Mandarin for a serious accident, then the Mandarin for who is to blame.
\end{quote}

\subsection*{You'll want to know: \zh{后果}}
\textbf{\zh{后果}} — \emph{hòuguǒ} — \textquotedblleft{}consequences; what comes of it\textquotedblright{}.

\begin{quote}
\emph{kāichē tài kuài, hòuguǒ huì hěn yánzhòng} — If you drive too fast, the consequences will be serious.
\end{quote}

\begin{cousinweb}
\textbf{\zh{后果}} (\emph{hòuguǒ}) is built from \zh{后} 'after' and \zh{果} 'fruit', both inherited from Middle Chinese, \emph{huwX} and \emph{kwaX} (Guangyun, in Baxter's transcription), which glosses the second 'the fruit of a tree' and 'resolute': the fruit that comes after.
\end{cousinweb}

\begin{grammarlens}[title={What you've built}]
One more word for telling what happened.
\end{grammarlens}

\subsection*{Your turn}
\begin{itemize}
\raggedright
  \item \emph{Say it:} \emph{hòuguǒ}
  \item \emph{Say it:} \emph{hòuguǒ}, once more
  \item \emph{Say it:} \emph{kāichē tài kuài, hòuguǒ huì hěn yánzhòng} — If you drive too fast, the consequences will be serious.
  \item \emph{Recall:} say the Mandarin for a serious accident, then the Mandarin for who is to blame, then say \emph{hòuguǒ} again
\end{itemize}

\begin{tcolorbox}[breakable,colback=teal!4,colframe=teal!35!black,title={Before you move on}]
What does \zh{后果} mean? (\textquotedblleft{}Consequences; what comes of it\textquotedblright{}.) Say it once more.
\end{tcolorbox}
\section[chǔlǐ]{chǔlǐ — to deal with}

\label{lesson:ZH-C297-chuli}



\begin{quote}
Before the new one: say the Mandarin for who is to blame, then the Mandarin for what comes of it.
\end{quote}

\subsection*{You'll want to know: chǔlǐ}
\emph{chǔlǐ} — \textquotedblleft{}to deal with\textquotedblright{}.

\begin{quote}
\emph{jǐngchá lái chǔlǐ le zhège shìgù} — The police came to deal with the accident.
\end{quote}

Its characters come later; for now it is a word for the ear.

Two third tones meet in this word: chǔ is said with a rising tone before lǐ, and the spelling keeps both marks.

\begin{cousinweb}
\emph{chǔlǐ} is built from \emph{chǔ} 'to settle' and \emph{lǐ} 'to put in order', both inherited from Middle Chinese, \emph{tsyhoX} and \emph{liX} (Guangyun, in Baxter's transcription), which glosses the first 'to dwell' and 'to settle' and the second 'to arrange' and 'reason'.
\end{cousinweb}

\begin{grammarlens}[title={What you've built}]
One more word for telling what happened.
\end{grammarlens}

\subsection*{Your turn}
\begin{itemize}
\raggedright
  \item \emph{Say it:} \emph{chǔlǐ}
  \item \emph{Say it:} \emph{chǔlǐ}, once more
  \item \emph{Say it:} \emph{jǐngchá lái chǔlǐ le zhège shìgù} — The police came to deal with the accident.
  \item \emph{Recall:} say the Mandarin for who is to blame, then the Mandarin for what comes of it, then say \emph{chǔlǐ} again
\end{itemize}

\begin{tcolorbox}[breakable,colback=teal!4,colframe=teal!35!black,title={Before you move on}]
What does chǔlǐ mean? (\textquotedblleft{}To deal with\textquotedblright{}.) Say it once more.
\end{tcolorbox}
